\chapter{Invariants and Exploitability}

\section{Purpose}

\begin{principlebox}
An exploit is often a counterexample to an intended invariant: a valid sequence of actions that reaches a state, transition, or history the protocol promised would not occur.
\end{principlebox}

Chapter 6 ended with a demand: a threat model must produce review work. This chapter teaches the most important unit of that work. An invariant is a property that should remain true across the states, transitions, or histories the system can actually reach. It is not a slogan such as ``funds are safe.'' It is a claim that can be challenged by a counterexample.

Many serious blockchain failures are not invalid transactions. They are valid transitions that the system accepts even though they violate protocol intent. A vault accepts a deposit that mints too few shares. A lending market accepts a borrow against a manipulated price. A bridge accepts the same message twice. A governance executor runs a payload before the intended delay. A withdrawal queue reaches a state where users have claims but no practical exit. The machine is doing what the implementation permits. The protocol is failing because the permitted behavior contradicts the intended security property.

The workflow for this chapter is:

\begin{codeblock}
protocol intent
    -> invariant
    -> counterexample
    -> exploitability analysis
    -> mitigation invariant
    -> tests, monitoring, and residual-risk decision
\end{codeblock}

The last phrase matters. A broken invariant is not automatically a catastrophic exploit. Exploitability depends on capital, timing, ordering, liquidity, permissions, reliability, profit, and non-profit impact. A reviewer must be able to say not only ``this property can fail,'' but also ``here is the sequence, here are the preconditions, here is what the attacker gains or breaks, and here is what would make the claim no longer true.''

\section{The Invariant Model}

\subsection{Reachable Properties}

For a state-machine view of a system, an invariant is usually written as:

\begin{codeblock}
for every reachable state S:
    P(S) == true
\end{codeblock}

The word \emph{reachable} does the work. A property that holds only in clean, intended, demo states is not a security invariant. The reviewer cares about states reachable through adversarial but valid transactions, odd timing, callbacks, direct token transfers, stale oracle data, governance changes, retries, pauses, and partial failures.

The proof idea is old and simple. Establish the property at initialization, then show that each allowed transition preserves it. Lamport's state-machine framing uses initial states and next-state relations to reason about invariance over reachable states.\footnote{\href{https://lamport.azurewebsites.net/pubs/state-machine.pdf}{Leslie Lamport, ``Computation and State Machines''.}} Smart-contract tools use the same basic idea in practical form: generate or reason over sequences of calls, then check whether a property still holds after those calls.\footnote{\href{https://www.getfoundry.sh/guides/invariant-testing}{Foundry, ``Invariant Testing''.}}

In security review, we often search in the opposite direction:

\begin{codeblock}
Can an attacker find a valid transition sequence
from an allowed initial state
to a state, transition, or history where P is false?
\end{codeblock}

That sequence is a counterexample. In a blockchain system, the counterexample is often already close to an exploit trace.

\subsection{State, Transition, and History Properties}

Not every important property is a simple predicate over one state. Some properties constrain a before/after relation. Some constrain event order. Some are economic assumptions. Some say that something must eventually become possible. Forcing every claim into one shape makes the specification weaker, not stronger.

\begin{table}[htbp]
\centering
\small
\renewcommand{\arraystretch}{1.18}
\begin{tabularx}{\textwidth}{@{}>{\RaggedRight\arraybackslash}p{0.22\textwidth}>{\RaggedRight\arraybackslash}p{0.28\textwidth}X@{}}
\toprule
Property type & Question & Example \\
\midrule
State invariant & What must be true now? & Recoverable assets must cover user claims under stated assumptions \\
Transition invariant & What must hold across one action? & A withdrawal may reduce a user's shares only by the amount redeemed \\
Relational invariant & Which records must agree? & Sum of account debts must agree with market debt after interest indexing \\
Temporal property & What must eventually or never happen? & An accepted withdrawal becomes claimable under documented conditions \\
Ordering property & What must happen before what? & A governance payload cannot execute before the timelock has elapsed \\
Economic invariant & What value relation must hold? & A healthy account cannot become undercollateralized through an accepted borrow \\
\bottomrule
\end{tabularx}
\caption{Useful invariants have the right shape for the property being claimed.}
\end{table}

Safety properties say bad things do not happen. Liveness properties say good things eventually happen. Alpern and Schneider give the classical safety/liveness distinction.\footnote{\href{https://www.cs.cornell.edu/fbs/publications/DefLiveness.pdf}{Bowen Alpern and Fred B. Schneider, ``Defining Liveness''.}} In blockchain security, theft, unauthorized minting, and bad debt are usually safety failures. Stuck withdrawals, unusable exits, unprocessable queues, and recovery paths that cannot be called are usually liveness failures. Both are security failures.

\subsection{Intended Invariants and Implementation Facts}

A true statement about the implementation can still be useless.

\begin{codeblock}
Weak:
    totalSupply is always equal to totalSupply.

Too local:
    storedTotalDeposits decreases only in withdraw().

Useful:
    The accounting variables that define user claims correspond to
    actual recoverable assets under the token, strategy, oracle, and
    withdrawal assumptions the protocol supports.
\end{codeblock}

The second statement may be true and still miss the bug. Token balances may change through direct transfers, rebases, fees, hooks, slashing, strategy losses, or cross-chain minting. Security review cares about intended invariants: the properties the protocol must satisfy for its security story to be true.

The gap between implemented facts and intended invariants is where deep bugs live.

\section{Writing Useful Invariants}

\subsection{From Wish to Property}

An invariant should be specific enough that a reviewer can imagine a counterexample.

\begin{codeblock}
Wish:
    Users should not lose funds.

Invariant:
    A successful withdrawal of s shares by user u must transfer no less
    than the assets owed for s shares under the protocol's current
    accounting rule, minus documented fees and bounded rounding.
\end{codeblock}

\begin{codeblock}
Wish:
    Only governance can upgrade the protocol.

Invariant:
    The proxy implementation address can change only through the
    governance executor after a passed proposal, elapsed timelock,
    matching payload hash, and unused execution record.
\end{codeblock}

\begin{codeblock}
Wish:
    The bridge is backed.

Invariant:
    For each asset, total destination-domain wrapped supply plus
    pending unlock obligations must not exceed source-domain locked
    balance plus explicitly authorized credit, under the bridge's
    finality and message-validity assumptions.
\end{codeblock}

The difference is reviewability. A wish cannot be disproved. A real invariant can be attacked.

\subsection{Quality Test}

Before using an invariant in review, testing, fuzzing, or formal verification, test the invariant itself.

\begin{table}[htbp]
\centering
\small
\renewcommand{\arraystretch}{1.18}
\begin{tabularx}{\textwidth}{@{}>{\RaggedRight\arraybackslash}p{0.24\textwidth}X@{}}
\toprule
Test & What the invariant must do \\
\midrule
Specific & Name the state, asset, role, domain, transition, or history being constrained \\
Falsifiable & Admit a concrete counterexample if it is false \\
Reachability-aware & Apply to adversarially reachable states, not only happy paths \\
Assumption-aware & Name oracle, token, finality, liquidity, governance, or operator assumptions \\
Quantified & Say whether it applies to all users, assets, markets, messages, or a subset \\
Measurable & Be checkable from code, state, traces, events, models, or explicit review evidence \\
Non-circular & Avoid treating an implementation variable as truth merely because it exists \\
Impact-linked & Matter if violated \\
\bottomrule
\end{tabularx}
\caption{A weak invariant makes the rest of the review weaker.}
\end{table}

Tools amplify good properties. They also amplify bad ones. Echidna can generate arbitrary transaction sequences and report a sequence that falsifies a property.\footnote{\href{https://secure-contracts.com/program-analysis/echidna/introduction/how-to-test-a-property.html}{Trail of Bits, ``Testing a Property with Echidna''.}} A formal prover can check preservation of a specification. Certora's invariant documentation describes establishment and preservation checks for smart contracts.\footnote{\href{https://docs.certora.com/en/latest/docs/cvl/invariants.html}{Certora Prover Documentation, ``Invariants''.}} But the tool does not know what the protocol should mean. The reviewer must choose the property.

\section{Invariant Families in Blockchain Systems}

\subsection{Conservation and Backing}

For systems that mint, burn, lock, unlock, wrap, deposit, withdraw, stake, swap, or bridge assets, start with the conservation question:

\begin{codeblock}
Can claims exceed backing?
\end{codeblock}

Examples:

\begin{codeblock}
Token:
    total issued supply must match authorized balances and mint/burn records.

Vault:
    recoverable assets must cover shares and pending withdrawal claims.

Bridge:
    wrapped supply and pending unlocks must be backed by locked source
    assets or explicitly authorized credit.

Reward distributor:
    claimed rewards + unclaimed rewards + remaining funded balance
    must not exceed funded rewards.
\end{codeblock}

Do not confuse exact conservation with bounded drift. Rounding, fees, strategy yield, bad debt, slashing, and asynchronous settlement may make equality wrong. But drift must be bounded, intentional, and assigned. A useful invariant names units, precision, decimal scaling, index scaling, and acceptable rounding.

\subsection{Authorization and Capability}

Authorization is not just a function modifier. It is a property of a state transition.

\begin{codeblock}
No transition may reduce user u's asset claim unless one is true:
    u authorized the transition,
    u's position satisfies the liquidation predicate,
    a documented fee was charged according to protocol rules,
    or an explicitly scoped emergency process executed.
\end{codeblock}

The invariant should bind authority to the protected object, action, verifier, chain or domain, nonce, and validity window. This applies to signatures, ownership checks, program-derived addresses, multisig modules, governance executors, session keys, bridge signer sets, and upgrade roles.

Negative authorization matters too. A guardian may pause deposits without being allowed to withdraw user funds. A liquidator may seize only collateral tied to an unhealthy debt. A strategist may move assets only into approved strategies. A role without negative boundaries is not a role; it is a vague trust assumption.

\subsection{Solvency, Prices, and Shares}

Solvency is a claim about recoverable assets versus enforceable claims under stated conditions.

\begin{codeblock}
recoverable assets
    >= immediately withdrawable claims
     + pending withdrawal claims
     + accrued obligations
     + recognized bad debt or loss socialization
\end{codeblock}

For lending systems, the property must include prices and operation classes:

\begin{codeblock}
Healthy-only operations must not leave an account below the
protocol's collateralization predicate under an accepted price snapshot.

Liquidation-only operations may execute only when the account satisfies
the documented liquidation predicate.
\end{codeblock}

The phrase ``accepted price snapshot'' hides real assumptions: oracle source, freshness, deviation bounds, sequencer availability, liquidity, decimals, and governance-updated risk parameters.

Vault shares are another recurring solvency surface. ERC-4626 standardizes tokenized vault interfaces such as \texttt{totalAssets}, \texttt{convertToShares}, and \texttt{convertToAssets}.\footnote{\href{https://eips.ethereum.org/EIPS/eip-4626}{Ethereum EIP-4626, ``Tokenized Vaults''.}} A share invariant should say what a share claims, what fees and rounding are allowed, what happens at zero or low supply, and whether direct donations affect pricing.

\subsection{Time, Exits, and Recovery}

Not every security failure is theft. A protocol can preserve accounting while trapping users.

\begin{codeblock}
If a withdrawal request is accepted at epoch e,
and required liquidity becomes available,
and no documented emergency state is active,
then the request must become claimable according to the queue rule,
without later requests bypassing it except under documented priority rules.
\end{codeblock}

Temporal properties need clocks and conditions. A timelock invariant should name proposal state, queued time, delay, payload hash, expiration, cancellation, and execution record. A bridge exit invariant should name burn, message finality, replay prevention, pause state, and unlock. A pause invariant should say which actions stop, which recovery actions remain available, and which new powers the pause authority does not gain.

Recovery controls need invariants because they intentionally bypass normal operation. A recovery action should move the system from a degraded state to a documented safe state without silently increasing user claims, reducing user claims arbitrarily, or granting permanent emergency authority.

\section{Counterexamples and Attack Hypotheses}

An invariant becomes useful when it creates an attack question.

\begin{codeblock}
Invariant:
    A borrower must not increase debt if the resulting account is
    undercollateralized under the accepted price snapshot.

Attack question:
    Can an attacker make the accepted price snapshot temporarily wrong,
    borrow against inflated collateral, and keep the debt after the
    price normalizes?
\end{codeblock}

\begin{codeblock}
Invariant:
    Only the owner of a position can withdraw collateral unless the
    position is liquidatable.

Attack question:
    Can an attacker make the system treat them as owner, make a healthy
    position appear liquidatable, or call a path that bypasses the check?
\end{codeblock}

The invariant is not only a test assertion. It is a generator of hypotheses.

\begin{figure}[htbp]
\centering
\begin{tikzpicture}[
  node distance=8mm,
  flow/.style={security node, text width=30mm, minimum height=8mm}
]
\node[flow] (intent) {Protocol intent};
\node[flow, right=of intent] (inv) {Invariant};
\node[flow, right=of inv] (trace) {Counter-\\example\\trace};
\node[flow, below=of trace] (exploit) {Exploitability};
\node[flow, left=of exploit] (mitigation) {Mitigation invariant};
\node[flow, left=of mitigation] (tests) {Tests and monitoring};

\draw[security arrow] (intent) -- (inv);
\draw[security arrow] (inv) -- (trace);
\draw[security arrow] (trace) -- (exploit);
\draw[security arrow] (exploit) -- (mitigation);
\draw[security arrow] (mitigation) -- (tests);
\end{tikzpicture}
\caption{Invariant review turns intent into a falsifiable trace and then into evidence.}
\end{figure}

The smallest counterexample is especially valuable. It removes noise. If an exploit requires ten steps, ask whether three of them are essential. If a sequence needs five contracts, ask which cross-contract assumption first failed. If the attack needs a particular deployment state, write it as a precondition rather than hiding it in prose.

\section{Exploitability}

\subsection{From Weakness to Attack}

A vulnerability is a weakness. An exploit is a feasible strategy that uses the weakness to achieve impact. The same invariant violation can be critical in one deployment and low severity in another because the exploitability conditions differ.

\begin{table}[htbp]
\centering
\small
\renewcommand{\arraystretch}{1.18}
\begin{tabularx}{\textwidth}{@{}>{\RaggedRight\arraybackslash}p{0.22\textwidth}X@{}}
\toprule
Factor & Questions \\
\midrule
Initial state & Must the vault be empty, an account unhealthy, a queue full, or a role misconfigured? \\
Capital & Must funds be owned, flash-borrowed, locked, staked, or exposed to market risk? \\
Timing and ordering & Must the attacker front-run, back-run, control ordering, wait for a stale price, or act before finality? \\
Liquidity & Is there enough depth to manipulate price, enter, exit, or realize profit? \\
Permissions & Does the path require public access, a role, compromised key, governance power, or operator collusion? \\
Reliability & Is the sequence deterministic, probabilistic, race-dependent, detectable, or repeatable? \\
Impact & Does it steal funds, create bad debt, freeze exits, seize control, grief users, or corrupt accounting? \\
\bottomrule
\end{tabularx}
\caption{Exploitability analysis scopes the counterexample to realistic attacker conditions.}
\end{table}

Use a worksheet for serious violations:

\begin{codeblock}
Invariant:
Violation:
Smallest counterexample:
Required initial state:
Required attacker capability:
Required victim, keeper, or operator behavior:
Required external dependency behavior:
Capital needed:
Timing or ordering needed:
Liquidity needed:
Profit or impact:
Repeatability:
Detection likelihood:
Mitigations:
Residual risk:
\end{codeblock}

``Requires a lot of capital'' is not analysis unless the amount, source, duration, and risk are specified. Capital that must be held across a challenge window is different from liquidity borrowed for one transaction. Voting power at a historical snapshot is different from tokens bought after the snapshot. A donation that is recoverable through redemption is different from a cost that is burned.

Exploitability should also define the boundary of the claim. A finding may be critical for new deployments but low risk for an already seeded instance. It may be unprofitable for a single market but profitable across many clones. It may require a malicious sequencer, a compromised admin key, or only a public call. The reviewer should state the strongest confirmed exploit, the weaker variants that remain plausible, and the deployment changes that would make severity increase. This prevents two common errors: dismissing a real invariant violation because one variant is blocked, and inflating severity without proving that the damaging path is reachable.

\subsection{Narrow Claims About Non-Exploitability}

Do not write:

\begin{codeblock}
This is not exploitable.
\end{codeblock}

Write:

\begin{codeblock}
Under the current deployment, the zero-share first-deposit attack
requires the vault to be empty. This vault has seeded dead shares,
a virtual offset, and deposits routed through a minimum-shares-out
check. A direct donation can change displayed share price, but the
specific victim-deposit theft path is not feasible unless those
conditions change.
\end{codeblock}

That statement is useful because it names the attack variant, deployment state, mitigations, remaining effect, and configuration sensitivity. ``Not exploitable'' without scope is usually false confidence.

Profit is also not the only impact. A cheap griefing attack that freezes a withdrawal queue can be severe even if the attacker cannot directly steal. A governance or upgrade takeover may have delayed impact. A liveness failure can convert a liquid claim into an indefinite promise. Severity should track impact, not only attacker revenue.

\section{Worked Example: Vault Donation and Share Inflation}

\subsection{Protocol Intent}

Consider a simplified tokenized vault:

\begin{codeblock}
Users deposit an ERC-20 asset and receive vault shares.
Users redeem shares for a proportional amount of vault assets.
The vault uses asset.balanceOf(vault) as total assets.
When totalShares == 0, initial deposits mint shares 1:1 with assets.
When totalShares > 0:
    sharesOut = assetsIn * totalShares / asset.balanceOf(vault)
rounded down.
\end{codeblock}

The intended security story is that shares represent proportional claims, depositors are not diluted unexpectedly, rounding is small, and anyone may deposit. Direct ERC-20 transfers into the vault cannot generally be prevented, so they must be modeled.

\subsection{Invariant}

Basic accounting is not enough:

\begin{codeblock}
totalShares == sum(userShares)
asset.balanceOf(vault) >= assetsRedeemableByAllShares
\end{codeblock}

Those statements may remain true during the attack. The deeper property is deposit fairness:

\begin{codeblock}
For any successful deposit of assets a,
the shares minted to the receiver must represent the receiver's
economic contribution, with loss bounded by documented fees,
slippage controls, and rounding policy.
\end{codeblock}

With adversarial context:

\begin{codeblock}
A third party must not be able to manipulate the deposit exchange rate
through direct donation so that a victim receives fewer shares than
the protocol's tolerance allows while the manipulator recovers the
donation through existing shares.
\end{codeblock}

OpenZeppelin's ERC-4626 documentation describes this class of inflation attack and the virtual-offset defense.\footnote{\href{https://docs.openzeppelin.com/contracts/5.x/erc4626}{OpenZeppelin Contracts documentation, ``ERC4626''.}}

\subsection{Counterexample Trace}

\begin{figure}[htbp]
\centering
\begin{tikzpicture}[
  node distance=7mm,
  trace/.style={security node, text width=31mm, minimum height=8mm}
]
\node[trace] (s0) {Empty or low-supply vault};
\node[trace, right=of s0] (s1) {Attacker deposits 1 unit};
\node[trace, right=of s1] (s2) {Attacker donates assets};
\node[trace, below=of s2] (s3) {Victim deposit rounds down};
\node[trace, left=of s3] (s4) {Attacker redeems};
\node[trace, left=of s4] (s5) {Victim loses claim};

\draw[security arrow] (s0) -- (s1);
\draw[security arrow] (s1) -- (s2);
\draw[security arrow] (s2) -- (s3);
\draw[security arrow] (s3) -- (s4);
\draw[security arrow] (s4) -- (s5);
\end{tikzpicture}
\caption{The attack is a valid transition sequence that violates deposit fairness.}
\end{figure}

The concrete trace:

\begin{codeblock}
1. Vault has zero or tiny share supply.
2. Attacker deposits 1 atomic unit and receives 1 share.
3. Attacker transfers D assets directly to the vault.
4. Victim deposits U assets without minimum-share protection.
5. sharesOut rounds to zero or an unfairly low value.
6. Attacker redeems the dominant share position.
7. Attacker captures a material part of the victim deposit.
\end{codeblock}

The broken invariant is not ``share price increased.'' Share price may legitimately increase through yield or donations. The broken invariant is that an attacker-controlled donation can make another user's deposit mint an economically unfair claim and can route the victim's loss back to the attacker.

\subsection{Exploitability}

The attack is strongest when the vault is empty or has very low share supply. The attacker needs enough assets to make the victim's share output round down to the target level. The donation may be mostly recoverable if the attacker is the dominant shareholder, so the gross capital requirement can overstate the net cost.

Timing matters. The attacker benefits from preceding the victim deposit, either by front-running, by initializing before public users arrive, or by exploiting a predictable first deposit. Liquidity may matter less than in an oracle manipulation because the attack uses the vault asset itself, but funding source still matters. If the deposit, donation, victim deposit, and redemption can be bundled or financed temporarily, the barrier is lower.

Reliability depends on deployment state and user path. If deposits enforce minimum shares out, the victim transaction can revert. If the vault is seeded with dead shares, virtual shares and assets, or enough initial liquidity, the attack may become uneconomic. If public users can deposit directly into unseeded vaults, the risk is high.

\subsection{Mitigation Invariants}

Mitigations are not magic words. Each one changes the property that must hold.

\begin{codeblock}
Minimum shares out:
    A deposit must revert if sharesOut < minSharesOut supplied by the caller.

Virtual offset:
    The initial exchange rate must be anchored so direct donation cannot
    make victim rounding loss exceed the documented bound at rational
    capital levels.

Seed liquidity:
    Public deposits must not open until initial share supply and assets
    bound rounding loss for expected deposit sizes.

Internal accounting:
    Direct token transfers must not affect deposit pricing unless the
    protocol explicitly treats donations as yield and bounds donation
    capture risk.

Monitoring:
    Alert when a low-supply vault receives a direct asset transfer near
    public deposit activity.
\end{codeblock}

The review is complete only when the mitigation has its own invariant, tests, and residual-risk statement.

\section{Invariant Catalogue}

For a serious protocol, maintain a compact invariant catalogue. It should be short enough to use and precise enough to guide testing, formal work, manual review, and monitoring.

\begin{codeblock}
Invariant ID:
Name:
Type:
    state / transition / temporal / economic / authorization / cross-domain
Statement:
Scope:
    contracts, assets, users, domains, roles
Assumptions:
    oracle, token behavior, finality, governance, operators, liquidity
Allowed exceptions:
Impact if violated:
Likely attack paths:
Tests:
    unit:
    stateful/fuzz:
    scenario:
    fork/simulation:
    formal/spec:
Monitoring:
Residual risk:
\end{codeblock}

Filled example:

\begin{codeblock}
Invariant ID:
    VAULT-ROUNDING-001

Name:
    Deposit rounding loss is bounded

Type:
    transition / economic

Statement:
    For any successful deposit of assets a, shares minted to receiver
    must satisfy the caller's minimum output and must not create
    attacker-capturable rounding loss above configured tolerance.

Scope:
    All vault deposit and mint paths, including router and direct calls.

Assumptions:
    Asset decimals are configured correctly.
    Direct donations are ignored by internal accounting or made uneconomic.
    Asset transfer behavior is supported and tested.

Impact if violated:
    User deposits can be partially or fully captured by existing shareholders.

Likely attack paths:
    First-deposit donation.
    Low-supply rounding manipulation.
    Fee-on-transfer mismatch.

Tests:
    Empty-vault first deposit.
    Stateful sequence: deposit, donate, victim deposit, redeem.
    Fork test against supported assets with real decimals.

Monitoring:
    Alert on direct transfers into low-supply vaults.
\end{codeblock}

This catalogue becomes the handoff from Part I's security model to later testing, simulation, formal verification, monitoring, and review-workflow chapters.

\section*{Review Questions}

\begin{enumerate}
\item Why is ``users should not lose funds'' not a useful invariant?
\item What is the difference between a property over all imaginable states and a property over reachable states?
\item Give one example each of a state invariant, transition invariant, temporal property, and economic invariant.
\item Why can a true implementation fact still fail to capture protocol intent?
\item What assumptions would you attach to a lending collateralization invariant?
\item Why is \texttt{totalSupply == sum(balances)} insufficient for a wrapped asset or vault share?
\item What must be specified before claiming that a bug is ``not exploitable''?
\item Why are liveness failures security failures even when no funds are immediately stolen?
\end{enumerate}

\section*{Applied Exercises}

\begin{enumerate}
\item For a vault that deposits assets into external strategies and processes withdrawals once per epoch, write an invariant catalogue with at least five invariants: conservation, share pricing, strategist authority, withdrawal ordering, and recovery after strategy failure. State assumptions for each.

\item A lending market accepts an AMM spot price as collateral input. Write the collateralization invariant, identify the hidden oracle and liquidity assumptions, then produce one counterexample trace and one exploitability worksheet.

\item A bridge locks ETH on chain A and mints wrapped ETH on chain B. Write a cross-domain conservation invariant that includes pending messages, finality delay, burns, unlocks, and replay prevention. Then write two attack hypotheses that would violate it.

\item Turn this counterexample into a short finding: an empty vault lets an attacker deposit 1 unit, donate 100,000 units directly, cause a victim deposit of 50,000 units to mint 0 shares, then redeem almost all assets. Include root cause, broken invariant, preconditions, impact, mitigation invariant, tests, and residual risk.
\end{enumerate}
